\chapter{Grandes volúmenes de datos}
\label{cap:bigdata}
\textit{Criterio 2 de la rúbrica (12\,\%): integra diversas fuentes; estructura, documenta, transforma y almacena
considerando calidad y arquitectura.}

\section{La arquitectura medallón: Bronze, Silver y Gold}
\begin{intuicion}
Piensa en una planta potabilizadora. El agua cruda (Bronze) se guarda tal como llega, para poder volver a ella. Se filtra
y se limpia (Silver). Y al final se embotella lista para beber (Gold). Si alguien duda del agua embotellada, se puede
revisar cada paso hasta la toma original.
\end{intuicion}

\begin{center}
\begin{tabularx}{\linewidth}{llY}
\encabezado \h{Capa} & \h{Carpeta} & \h{Contenido y regla} \\
\textbf{Bronze} & \codigo{datos/bronze/} & Los archivos tal como los publica la fuente (zip, Excel, JSON, CSV, texto). Nunca se modifican; cada uno tiene su huella SHA-256. \\
\textbf{Silver} & \codigo{datos/silver/} & Una tabla limpia por fuente, en Parquet, con nombres en español, tipos correctos y banderas de calidad (\codigo{_flag}). La escribe Spark. \\
\textbf{Gold} & \codigo{datos/gold/} & Salidas de los modelos y el modelo estrella; lo que usan la página y los documentos. \\
\bottomrule
\end{tabularx}
\end{center}

\textbf{Por qué así.} Si una cifra de Gold parece rara, se puede repetir el camino desde el archivo crudo, que no se tocó.
\textbf{Alternativa descartada}: limpiar mientras se descarga. Así Bronze dejaría de ser idéntico a la fuente y ya no
habría contra qué comparar.

\subsection{La huella SHA-256}
\begin{intuicion}
Una huella SHA-256 es como la huella digital de un archivo: una cadena de 64 caracteres que se calcula con todo su
contenido. Si cambia un solo byte del archivo, la huella cambia por completo. No se puede ``fabricar'' otro archivo con la
misma huella.
\end{intuicion}
La prueba \codigo{tests/test_ingesta.py} vuelve a calcular las 353 huellas y falla si alguna no coincide con el
manifiesto. Así se demuestra que Bronze no se alteró.

\section{Spark: qué es y qué hace exactamente aquí}
\begin{intuicion}
Limpiar una lista de 6 millones de negocios es como corregir 6 millones de exámenes. Una persona sola tardaría mucho;
doce personas, cada una con un montón, terminan mucho antes. Spark reparte la tabla en montones (\emph{particiones}) y
pone a trabajar a todos los núcleos de la computadora a la vez.
\end{intuicion}

\subsection{Cinco ideas que debes poder explicar}
\begin{enumerate}
\item \textbf{Driver y ejecutores.} El \emph{driver} es el programa que planea el trabajo; los \emph{ejecutores} lo
  hacen. Aquí Spark corre en \textbf{modo local} (\codigo{local[*]}): el driver y los ejecutores viven en la misma
  computadora y usan sus 12 núcleos, con 4 GB de memoria. Es el mismo programa que correría en un clúster de muchas
  computadoras; solo cambia la dirección del ``maestro''.
\item \textbf{DataFrame y particiones.} Un DataFrame de Spark es una tabla repartida en particiones. Cada partición la
  procesa un núcleo. Se programa parecido a pandas (\codigo{select}, \codigo{withColumn}, \codigo{groupBy}).
\item \textbf{Evaluación perezosa.} Las \emph{transformaciones} (\codigo{withColumn}, \codigo{filter}) no se ejecutan al
  escribirlas: Spark arma un plan (un grafo acíclico de pasos) y lo optimiza completo. Solo cuando se pide una
  \emph{acción} (\codigo{count}, \codigo{write}) se ejecuta. Así puede, por ejemplo, filtrar antes de leer columnas que
  no se van a usar.
\item \textbf{Barajeo (shuffle).} Algunas operaciones necesitan juntar renglones de distintas particiones: quitar
  duplicados por \codigo{id} obliga a mandar todos los renglones con el mismo \codigo{id} al mismo núcleo. Es la parte
  cara del trabajo y por eso se hace una sola vez.
\item \textbf{Parquet particionado.} Spark escribe cada tabla en formato \textbf{Parquet}: guarda por columnas (leer una
  columna no obliga a leer las demás) y comprime. Además la parte en carpetas por una columna, por ejemplo
  \codigo{cve_ent=23/} para Quintana Roo. Leer solo Quintana Roo no obliga a leer los otros 31 estados
  (\emph{partition pruning}).
\end{enumerate}

\subsection{El trabajo más grande: los 6.1 millones de negocios del DENUE}
\begin{lstlisting}[language=Python]
crudo = (spark.read.option("header", True).option("encoding", "ISO-8859-1")
         .csv("temporal/*.csv"))                               # 33 CSV como UNA sola tabla
silver = (crudo.select(*COLUMNAS)                               # sin datos personales
          .dropDuplicates(["id"])                               # un negocio, una fila  (barajeo)
          .withColumn("scian_3", F.substring("codigo_act", 1, 3))
          .withColumn("categoria_turistica", categoria)         # 721 alojamiento, 722 alimentos...
          .withColumn("coordenadas_flag",
                      ~(lat.between(14, 33) & lon.between(-119, -86)) | lat.isNull())
          .withColumn("es_qroo", F.col("cve_ent") == "23"))
silver.write.mode("overwrite").partitionBy("cve_ent").parquet("datos/silver/denue")
\end{lstlisting}

Qué pasa en cada línea:
\begin{itemize}
\item \textbf{Lectura}: los 33 archivos (el Estado de México viene en dos partes) se leen como una sola tabla, con la
  codificación ISO-8859-1 que usa el INEGI.
\item \textbf{Privacidad}: se seleccionan solo las columnas necesarias; la razón social, el teléfono y el correo nunca
  entran a Silver.
\item \textbf{Duplicados}: \codigo{dropDuplicates(["id"])}. Resultado: 0 duplicados.
\item \textbf{Categoría turística}: con los primeros dígitos del código SCIAN (721 alojamiento, 722 alimentos y bebidas,
  5615 agencias de viajes, 487 transporte turístico, 712 museos y sitios, 713 esparcimiento).
\item \textbf{Bandera de calidad}: las coordenadas fuera de México se \emph{marcan}, no se borran (resultado: ninguna fuera).
\item \textbf{Escritura}: Parquet partido por estado. Tarda cerca de minuto y medio.
\end{itemize}

\subsection{Dónde se usa Spark y dónde no}
\begin{itemize}
\item \textbf{Sí}: las 12 limpiezas \codigo{silver_*.py}, la ventana deslizante de la Torre (\codigo{envivo/senales.py})
  y la reproducción con Structured Streaming (\codigo{envivo/torre.py}).
\item \textbf{No}: Spark no lee Excel, así que pandas abre los Excel de DataTur e INAH y Spark hace la unión, la
  deduplicación y la escritura. Los modelos (scikit-learn, statsmodels, PuLP) trabajan sobre tablas de Gold que ya caben
  en memoria: no tendría sentido repartirlas.
\end{itemize}

\begin{teprofe}
\emph{``¿De verdad se necesitaba Spark?''} Con honestidad: 8.1 millones de registros caben en una laptop y pandas podría
con casi todo. Spark se usó porque la rúbrica pide una arquitectura de grandes volúmenes, porque el DENUE nacional y el
clima por hora ya son pesados, y porque el mismo código escala sin cambios a un clúster si mañana entran más años o más
estados. El valor está en la arquitectura y el rastreo, no en el tamaño.
\end{teprofe}

\section{Reglas de limpieza: cero no es lo mismo que hueco}
La regla más importante de Silver: \textbf{un cero físicamente imposible es un hueco}, y un cero posible se conserva.

\begin{center}
\begin{tabularx}{\linewidth}{lYl}
\encabezado \h{Caso} & \h{Por qué} & \h{Regla} \\
Ocupación con 0 cuartos disponibles & No se pueden ocupar cuartos que no existen & Hueco \\
Los 4 aeropuertos en 0 todo 2025 & Cancún no pudo recibir cero pasajeros (su ocupación semanal promedió 72.2\,\%): la fuente dejó de publicar & Hueco \\
Afluencia y derrama en 0 & Un destino no recibe cero turistas en un mes & Hueco \\
Cruceristas en 0 en 2020 & Real: los puertos cerraron por la pandemia & Se conserva \\
Zona arqueológica con 0 visitantes & Suele ser un cierre por obras & Se conserva y se marca \\
Feriado sin tipo de cambio & No hubo cotización ese día & Vacío; no se copia el día anterior \\
\bottomrule
\end{tabularx}
\end{center}

\subsection{Ocupación: sumar antes de dividir}
La ocupación de un periodo \textbf{nunca} se calcula promediando porcentajes:
\[
O=\frac{\sum_{\text{meses}} \text{cuartos-noche ocupados}}{\sum_{\text{meses}} \text{cuartos-noche disponibles}}\times 100 .
\]
\textbf{Fundamento.} Un promedio simple de porcentajes le da el mismo peso a un mes con 50,000 cuartos-noche que a uno con
80,000. El cociente de sumas es el promedio \emph{ponderado} correcto: cada cuarto-noche cuenta una vez.

\begin{amano}[ · ocupación de Chetumal en 2024]
\[
O=\frac{458{,}696}{791{,}016}\times 100 = 57.99 \approx 58.0\,\%.
\]
De cada 10 cuartos, 4 se quedaron vacíos: capacidad ociosa, justo lo que la campaña puede aprovechar.
\end{amano}

\textbf{Un error real que lo demuestra.} SITUR-Q publica tres renglones por mes para la ocupación (disponibles, ocupados
y porcentaje). La primera versión del panel del Radar los sumaba, y la ``ocupación'' llegaba a 2,560,068. Ahora se
recalcula ocupados ÷ disponibles, y la prueba \codigo{test_ocupacion_no_suma_filas} lo vigila.

\section{Cuando dos fuentes cuentan lo mismo (conciliación)}
El incidente de Big Data del Problema Prototípico pregunta qué hacer con ``fuentes que no coinciden''. El proyecto tiene
dos casos medidos:
\[
\Delta\% = \left(\frac{S}{D}-1\right)\times 100, \qquad S=\text{SITUR-Q (estatal)},\quad D=\text{DataTur (federal)} .
\]

\begin{amano}[ · cruceristas en 2025]
Cozumel: \(\left(\dfrac{4{,}915{,}242}{4{,}724{,}255}-1\right)\times100 = (1.0404-1)\times100 = +4.04\,\%\).

Mahahual: \(\left(\dfrac{2{,}641{,}695}{2{,}379{,}422}-1\right)\times100 = +11.02\,\%\). En Mahahual SITUR-Q cuenta
siempre más: entre 11\,\% y 35\,\% según el año.
\end{amano}

En la ocupación pasa lo mismo: DataTur marca menos que SITUR-Q (Cancún $-1.6$ puntos, Cozumel $-14.5$, Isla Mujeres
$-19.3$), aunque las dos se mueven juntas (correlación entre 0.68 y 0.97).

\textbf{La regla del proyecto} tiene tres partes:
\begin{itemize}
\item cada lugar usa \textbf{una sola fuente en toda su historia}, así sus cambios en el tiempo son reales;
\item la diferencia se \textbf{declara} y no se ``ajusta'';
\item si la diferencia va en sentido conservador, la conclusión no cambia. Aquí el norte se ve \emph{menos} lleno con
  DataTur, así que ``el sur tiene espacio'' se sostiene aún con más fuerza.
\end{itemize}

\section{El almacén DuckDB y el modelo estrella}
\textbf{DuckDB} es una base de datos analítica que vive en un solo archivo (\codigo{datos/gold/torre.duckdb}) y lee los
Parquet de Silver y Gold \textbf{sin copiarlos}. Responde en milisegundos sin levantar Spark. Tiene una vista por cada
tabla (14 de Silver y 39 de Gold) y un \textbf{modelo estrella}.

\textbf{Qué es un modelo estrella} (Kimball): una tabla de \emph{hechos} en el centro, con las medidas, y tablas de
\emph{dimensiones} alrededor, con las descripciones. Su ventaja es que cualquier pregunta se contesta con el mismo patrón:
hechos unidos con dimensiones, filtrados y agrupados.

\begin{center}
\begin{tabularx}{\linewidth}{llY}
\encabezado \h{Tabla} & \h{Tipo} & \h{Contenido} \\
\codigo{dim_lugar} & Dimensión & 15 lugares con su papel: 5 promovidos, 3 de referencia y 7 de comparación \\
\codigo{dim_tiempo} & Dimensión & Un renglón por mes, de 2012 a 2027 \\
\codigo{hechos_mes} & Hechos & 3,287 renglones: lugar $\times$ mes $\times$ variable, con su fuente \\
\bottomrule
\end{tabularx}
\end{center}

\textbf{El grano} (qué representa un renglón de hechos) es ``una medida de un lugar en un mes''. Un hueco \textbf{no tiene
renglón}: nunca se guarda un cero inventado. Ejemplo de consulta:
\begin{lstlisting}[language=SQL]
SELECT anio, SUM(valor) FROM hechos_mes JOIN dim_tiempo USING (periodo)
WHERE lugar = 'Chetumal' AND variable = 'llegadas_extranjeros_avion'
GROUP BY anio;      -- 2025: 250 extranjeros (Cancún: 9,408,423)
\end{lstlisting}

\textbf{El diccionario de datos} (\codigo{docs/datos/DICCIONARIO.md}) se \textbf{genera solo} desde el almacén: 56 tablas,
cada columna con su tipo, su porcentaje de vacíos, un ejemplo real y su significado. Si una columna de Silver no tiene
significado escrito, el programa se detiene. Así el diccionario no puede quedar desactualizado (en la revisión del 2 de
octubre esa regla detectó 12 tablas sin descripción).

\textbf{Alternativa descartada}: PostgreSQL. Exige instalar y levantar un servidor, y el proyecto debe correr sin
internet en cualquier computadora.

\section{Streaming: la Torre en vivo con Spark Structured Streaming}
\label{sec:streaming}
\begin{intuicion}
Una banda transportadora: cada semana llega una caja con los datos de esa semana. La Torre abre una caja a la vez, en el
orden en que llegaron, decide qué hacer con la campaña y pasa a la siguiente.
\end{intuicion}

\textbf{Cómo funciona} (\codigo{envivo/torre.py}):
\begin{enumerate}
\item Cada una de las 239 semanas reales (enero de 2022 a julio de 2026) se escribe como un archivo JSON en una carpeta,
  con hora de llegada creciente.
\item Spark \textbf{escucha} la carpeta con \codigo{readStream}. La opción \codigo{maxFilesPerTrigger = 1} hace que cada
  \emph{micro-lote} sea exactamente un archivo, es decir, una semana. \codigo{trigger(availableNow=True)} procesa todo lo
  que hay y se detiene.
\item En cada micro-lote, la función \codigo{foreachBatch} le pasa la semana al motor de reglas, que decide qué anuncio se
  enciende o se pausa y lleva la cuenta del dinero guardado.
\item Un \textbf{punto de control} (\emph{checkpoint}) registra qué archivos ya se procesaron: si el programa se cae, no
  procesa dos veces la misma semana.
\item Al final se comprueba que llegaron \textbf{239 lotes, en orden}.
\end{enumerate}

\textbf{Por qué así.} El incidente de Big Data pide una arquitectura de \emph{baja latencia} para ajustar la campaña
\emph{durante} su ejecución. Esta es exactamente la arquitectura que recibiría datos reales cada semana: una carpeta a la
que llega un archivo nuevo. Cambiar la fuente por una real no cambia el motor de reglas. En la página se muestra como
``reproducción de datos históricos reales'', porque GitHub Pages no puede transmitir en vivo.

\subsection{La ventana deslizante}
El promedio de ocupación de las últimas cuatro semanas lo calcula Spark con una \emph{ventana} sobre cada centro:
\[
\bar o_w=\frac14\sum_{k=0}^{3}o_{w-k}.
\]
La ventana ``se desliza'': cada semana entra la nueva y sale la de hace cuatro. Suaviza una semana atípica sin perder la
tendencia reciente.

\section{Privacidad}
No se guarda ningún dato personal. El DENUE entra sin razón social, teléfono ni correo; las reseñas de Rest-Mex no tienen
identificadores; las nacionalidades y los visitantes vienen agregados por mes. Ninguna tabla permite identificar a una
persona.
